% ============================================================================
% METHODS BRIEF: POLYPHARMACY RISK ASSESSMENT AND ALTERNATIVE THERAPY SUGGESTION
% Publication Format - Paragraph Style for Paper
% Based on Actual Knowledge Graph Data
% ============================================================================

\documentclass[11pt,a4paper]{article}
\usepackage[utf8]{inputenc}
\usepackage{amsmath,amssymb}
\usepackage{booktabs}
\usepackage{geometry}
\usepackage{graphicx}
\usepackage{hyperref}
\usepackage{float}
\usepackage{xcolor}
\geometry{margin=1in}

\hypersetup{
    colorlinks=true,
    linkcolor=blue,
    filecolor=magenta,
    urlcolor=cyan,
}

\title{\textbf{Polypharmacy Risk Assessment and Alternative Therapy Suggestion}\\
\large Methods Brief for Publication}
\author{Drug-Drug Interaction Knowledge Graph Project}
\date{}

\begin{document}

\maketitle

\section*{Abstract}

We present an integrated framework for polypharmacy risk assessment and therapeutic alternative recommendation built upon a comprehensive drug-drug interaction knowledge graph. The Polypharmacy Risk Index (PRI) quantifies individual drug risk using network-based metrics derived from 759,774 drug-drug interactions, while the Multi-Objective Drug Recommendation System identifies safer therapeutic substitutions. Applied to a high-risk cardiovascular regimen, the system demonstrated the ability to eliminate major interactions through evidence-based drug substitution.

\tableofcontents
\newpage

% ############################################################################
% PART I: METHODS
% ############################################################################

\part{Methods}

% ============================================================================
\section{Polypharmacy Risk Assessment}
% ============================================================================

\subsection{Polypharmacy Risk Index (PRI)}

We developed the Polypharmacy Risk Index (PRI) to quantify individual drug risk within multi-drug regimens. The PRI integrates four network-based metrics weighted according to their clinical relevance, combining degree centrality ($C_{\text{degree}}$), weighted degree centrality ($C_{\text{weighted}}$), betweenness centrality ($C_{\text{betweenness}}$), and severity profile score ($S$) into a composite risk measure:

\begin{equation}
\text{PRI}(d) = 0.25 \cdot C_{\text{degree}}(d) + 0.30 \cdot C_{\text{weighted}}(d) + 0.20 \cdot C_{\text{betweenness}}(d) + 0.25 \cdot S(d)
\label{eq:pri}
\end{equation}

\subsection{Component Metrics}

Degree centrality $C_{\text{degree}}(d)$ captures the interaction burden as the number of drugs interacting with drug $d$, normalized by the maximum possible degree in the network (weight: 0.25). This metric provides a direct measure of how many potential interaction partners a drug possesses. In our knowledge graph, drugs range from having zero interactions to over 4,700 documented interaction partners. \textit{Rationale:} The 0.25 weight reflects the fundamental importance of interaction count as a baseline risk indicator, while avoiding over-emphasis on quantity alone without considering severity.

Weighted degree centrality $C_{\text{weighted}}(d)$ extends the simple degree measure by incorporating severity information, representing the severity-weighted sum of all interactions (weight: 0.30). This gives greater importance to contraindicated and major interactions compared to moderate and minor ones, reflecting their clinical significance. \textit{Rationale:} This component receives the highest weight because it directly quantifies clinically significant interaction burden by combining both the number of interactions and their severity---the most actionable information for clinical decision-making.

Betweenness centrality $C_{\text{betweenness}}(d)$ quantifies the drug's role in risk propagation pathways across the network (weight: 0.20). Drugs with high betweenness centrality serve as ``bridges'' between interacting clusters and may facilitate cascade effects when multiple medications are combined. \textit{Rationale:} The lower weight acknowledges that while network topology provides valuable systems-level insight into potential cascade effects, direct interaction measures are more immediately relevant to individual patient risk assessment.

The severity profile score $S(d)$ measures the proportion of a drug's interactions that are classified as contraindicated or major severity (weight: 0.25). This metric identifies drugs that, while potentially having fewer total interactions, are predominantly involved in high-risk combinations. \textit{Rationale:} Equal weighting with degree centrality ensures that drugs with fewer but predominantly dangerous interactions are appropriately flagged, preventing false reassurance from low interaction counts when those interactions carry severe consequences.

\subsection{Risk Classification}

We defined heuristic PRI thresholds to stratify drugs into three risk categories based on the normalized 0-1 score range. Drugs with PRI scores exceeding 0.5 are classified as High Risk, flagged by the system for clinical review and consideration of therapeutic alternatives. PRI scores between 0.3 and 0.5 indicate Medium Risk, flagged for closer monitoring and potential medication review. Scores below 0.3 represent Lower Risk, where standard monitoring protocols may be appropriate (Table~\ref{tab:pri_thresholds}). These thresholds provide a practical framework for prioritizing clinical attention, though individual institutional protocols may adjust cutoffs based on local practice patterns.

\begin{table}[H]
\centering
\caption{PRI Risk Classification Thresholds}
\label{tab:pri_thresholds}
\begin{tabular}{lll}
\toprule
\textbf{PRI Score} & \textbf{Risk Level} & \textbf{Clinical Action} \\
\midrule
$> 0.5$ & High Risk & Flagged for clinical review \\
$0.3 - 0.5$ & Medium Risk & Flagged for closer monitoring \\
$< 0.3$ & Lower Risk & Standard monitoring protocols \\
\bottomrule
\end{tabular}
\end{table}

\subsection{Severity Weighting}

The severity weights employed in the weighted degree calculation correspond to the four severity categories from DrugBank's interaction classification. The weights are scaled to provide greater differentiation between severity levels. Contraindicated interactions receive the highest weight of 10.0, reflecting combinations that should generally be avoided. Major interactions are assigned a weight of 7.0, indicating significant clinical concern requiring intervention or close monitoring. Moderate interactions receive a weight of 4.0, suggesting the need for awareness and potential dose adjustments. Minor interactions are weighted at 1.0, representing clinically less significant but still documented interactions (Table~\ref{tab:severity_weights}).

\begin{table}[H]
\centering
\caption{Severity Weights for PRI Calculation}
\label{tab:severity_weights}
\begin{tabular}{lr}
\toprule
\textbf{Severity Level} & \textbf{Weight} \\
\midrule
Contraindicated & 10.0 \\
Major & 7.0 \\
Moderate & 4.0 \\
Minor & 1.0 \\
\bottomrule
\end{tabular}
\end{table}

% ============================================================================
\section{Alternative Therapy Suggestion}
% ============================================================================

\subsection{Multi-Objective Drug Recommendation System}

The recommendation system employs a multi-objective optimization framework to identify therapeutic alternatives that balance clinical efficacy with safety. For each candidate alternative drug $d_{\text{alt}}$, we compute a recommendation score that integrates therapeutic similarity, safety improvement, and overall risk reduction:

\begin{equation}
\text{RecScore}(d_{\text{alt}}) = 0.40 \cdot T(d_{\text{alt}}) + 0.35 \cdot \text{SAF}(d_{\text{alt}}) + 0.25 \cdot R(d_{\text{alt}})
\label{eq:recscore}
\end{equation}

\subsection{Therapeutic Similarity Score}

The therapeutic similarity score $T(d_{\text{alt}})$, weighted at 0.40, evaluates whether the alternative can serve the same clinical purpose as the original drug. This component ensures that recommended substitutions maintain therapeutic equivalence while potentially offering improved safety profiles.

Therapeutic similarity is computed as a weighted combination of four sub-scores. ATC code matching (weight 0.50) assesses pharmacological class equivalence based on the Anatomical Therapeutic Chemical classification system--drugs sharing the same ATC class are considered therapeutically interchangeable. Disease indication overlap (weight 0.25) confirms therapeutic applicability by comparing approved indications using Jaccard similarity. Shared protein targets (weight 0.15) quantifies similar mechanisms of action through common molecular targets. Metabolic pathway similarity (weight 0.10) evaluates comparable pharmacokinetic profiles:

\begin{equation}
T(d_{\text{alt}}) = 0.50 \cdot \text{ATC}(d_{\text{alt}}) + 0.25 \cdot \text{Disease}(d_{\text{alt}}) + 0.15 \cdot \text{Target}(d_{\text{alt}}) + 0.10 \cdot \text{Pathway}(d_{\text{alt}})
\end{equation}

\subsection{Safety Improvement Score}

The safety improvement score $\text{SAF}(d_{\text{alt}})$, weighted at 0.35, quantifies the reduction in DDI risk when the alternative is used with other drugs in the current regimen. This component is critical for achieving the primary goal of reducing adverse event potential:

\begin{equation}
\text{SAF}(d_{\text{alt}}) = \max\left(0, 1 - r_{\text{ddi}} - 0.3 \cdot n_{\text{contra}}\right)
\end{equation}

where $r_{\text{ddi}}$ is the normalized DDI risk score with the regimen and $n_{\text{contra}}$ is the count of contraindicated interactions introduced by the alternative. Higher scores indicate safer alternatives with fewer severe interactions.

\subsection{Risk Reduction Score}

The risk reduction score $R(d_{\text{alt}})$, weighted at 0.25, measures the net change in severe interactions and PRI score improvement when substituting the original drug with the alternative:

\begin{equation}
R(d_{\text{alt}}) = \frac{\Delta_{\text{contraindicated}} + \Delta_{\text{major}} + \alpha \cdot \Delta_{\text{PRI}}}{\text{max\_possible\_reduction}}
\end{equation}

where $\alpha$ is a scaling factor for PRI improvement contribution. This component captures the overall improvement in regimen safety profile, accounting for both categorical changes in interaction severity and continuous PRI improvements.

\subsection{Weight Selection Rationale}

The component weights were selected to prioritize clinical efficacy while balancing safety considerations. Therapeutic similarity receives the highest weight (0.40) to ensure that clinical efficacy is maintained as the primary consideration. Safety improvement is weighted at 0.35, reflecting its critical importance in reducing adverse events. Risk reduction receives a weight of 0.25 as a supporting metric for overall regimen optimization (Table~\ref{tab:weight_rationale}).

\begin{table}[H]
\centering
\caption{Weight Selection Criteria for Recommendation Score}
\label{tab:weight_rationale}
\begin{tabular}{lcp{8cm}}
\toprule
\textbf{Component} & \textbf{Weight} & \textbf{Rationale} \\
\midrule
Therapeutic Similarity & 0.40 & Highest priority to ensure clinical efficacy is maintained \\
Safety Improvement & 0.35 & Critical for reducing adverse events \\
Risk Reduction & 0.25 & Supporting metric for overall regimen optimization \\
\bottomrule
\end{tabular}
\end{table}

\newpage

% ############################################################################
% PART II: RESULTS
% ############################################################################

\part{Results}

% ============================================================================
\section{Knowledge Graph Statistics}
% ============================================================================

The DDI knowledge graph comprises 45,655 nodes distributed across six types: Drugs (4,313; 9.4\%), Proteins (3,176; 7.0\%), Side Effects (5,548; 12.2\%), Diseases (3,041; 6.7\%), Pathways (25,958; 56.9\%), and Categories (3,619; 7.9\%). The graph contains 1,211,674 edges across six relationship types, with drug-drug interactions comprising the largest portion (Table~\ref{tab:kg_stats}).

\begin{table}[H]
\centering
\caption{Knowledge Graph Edge Statistics}
\label{tab:kg_stats}
\begin{tabular}{lrrr}
\toprule
\textbf{Edge Type} & \textbf{Count} & \textbf{Percentage} & \textbf{Source} \\
\midrule
Drug-Drug Interaction & 759,774 & 62.7\% & DrugBank \\
Drug-Side Effect & 265,238 & 21.9\% & SIDER \\
Drug-Category & 70,618 & 5.8\% & DrugBank \\
Drug-Disease & 63,278 & 5.2\% & CTD \\
Drug-Pathway & 31,207 & 2.6\% & DrugBank \\
Drug-Protein & 21,559 & 1.8\% & DrugBank \\
\midrule
\textbf{Total} & \textbf{1,211,674} & \textbf{100.0\%} & -- \\
\bottomrule
\end{tabular}
\end{table}

% ============================================================================
\section{DDI Severity Distribution}
% ============================================================================

The knowledge graph contains 759,774 drug-drug interactions with severity classifications. The severity distribution shows that moderate interactions comprise the majority (464,655; 61.2\%), followed by minor interactions (211,991; 27.9\%), major interactions (83,102; 10.9\%), and contraindicated combinations (26; 0.003\%). This distribution reflects the clinical reality where severe interactions represent a smaller but clinically critical subset of all potential drug combinations (Table~\ref{tab:severity_distribution}).

\begin{table}[H]
\centering
\caption{DDI Severity Distribution (n = 759,774)}
\label{tab:severity_distribution}
\begin{tabular}{lrr}
\toprule
\textbf{Severity} & \textbf{Count} & \textbf{Percentage} \\
\midrule
Moderate & 464,655 & 61.2\% \\
Minor & 211,991 & 27.9\% \\
Major & 83,102 & 10.9\% \\
Contraindicated & 26 & 0.003\% \\
\midrule
\textbf{Total} & \textbf{759,774} & \textbf{100.0\%} \\
\bottomrule
\end{tabular}
\end{table}

% ============================================================================
\section{High-Risk Drugs by Interaction Count}
% ============================================================================

Analysis of the knowledge graph identified drugs with the highest interaction counts and severe interaction proportions. Quinidine leads with 4,782 total interactions (258 severe, 5.4\%), followed by Procaine with 4,719 interactions (786 severe, 16.7\%), and Lidocaine with 4,208 interactions (626 severe, 14.9\%). Notably, Warfarin, a commonly used anticoagulant, has 4,145 total interactions with 298 classified as severe (7.2\%), making it a frequent candidate for therapeutic optimization (Table~\ref{tab:high_risk_drugs}).

\begin{table}[H]
\centering
\caption{Top 10 Drugs by Total Interaction Count}
\label{tab:high_risk_drugs}
\begin{tabular}{llccc}
\toprule
\textbf{Drug} & \textbf{Total} & \textbf{Severe} & \textbf{Severe \%} & \textbf{Weighted} \\
\midrule
Quinidine & 4,782 & 258 & 5.4\% & 9,822 \\
Procaine & 4,719 & 786 & 16.7\% & 9,864 \\
Lidocaine & 4,208 & 626 & 14.9\% & 8,208 \\
Verapamil & 4,208 & 384 & 9.1\% & 6,840 \\
Propranolol & 4,184 & 208 & 5.0\% & 6,764 \\
Warfarin & 4,145 & 298 & 7.2\% & 8,586 \\
Isradipine & 4,142 & 162 & 3.9\% & 6,380 \\
Pindolol & 4,048 & 722 & 17.8\% & 7,244 \\
Indomethacin & 4,045 & 778 & 19.2\% & 7,536 \\
Nicardipine & 4,044 & 402 & 9.9\% & 6,766 \\
\bottomrule
\end{tabular}
\end{table}

% ============================================================================
\section{Case Study: High-Risk Cardiovascular Polypharmacy}
% ============================================================================

\subsection{Original Regimen Analysis}

We validated the methodology using a clinically realistic 5-drug cardiovascular regimen comprising Warfarin (anticoagulant), Amiodarone (antiarrhythmic), Digoxin (cardiac glycoside), Quinidine (antiarrhythmic), and Propranolol (beta-blocker). This combination represents a challenging but not uncommon clinical scenario in patients with complex cardiovascular conditions requiring multiple therapeutic agents (Table~\ref{tab:original_regimen}).

\begin{table}[H]
\centering
\caption{Original High-Risk Cardiovascular Regimen}
\label{tab:original_regimen}
\begin{tabular}{ll}
\toprule
\textbf{Drug} & \textbf{Therapeutic Class} \\
\midrule
Warfarin & Anticoagulant \\
Amiodarone & Antiarrhythmic \\
Digoxin & Cardiac Glycoside \\
Quinidine & Antiarrhythmic \\
Propranolol & Beta-blocker \\
\bottomrule
\end{tabular}
\end{table}

Analysis of this regimen in the knowledge graph revealed 10 unique drug-drug interaction pairs. One interaction was classified as Major: Warfarin-Amiodarone, which increases the risk or severity of bleeding when Amiodarone is combined with Warfarin. The remaining nine interactions were classified as Moderate, including pharmacokinetic interactions affecting serum concentrations and pharmacodynamic interactions affecting therapeutic efficacy (Table~\ref{tab:case_ddis}).

\begin{table}[H]
\centering
\caption{Drug-Drug Interactions in Cardiovascular Regimen}
\label{tab:case_ddis}
\begin{tabular}{lll}
\toprule
\textbf{Drug Pair} & \textbf{Severity} & \textbf{Mechanism} \\
\midrule
Warfarin + Amiodarone & Major & Increased bleeding risk \\
Warfarin + Digoxin & Moderate & Decreased Warfarin excretion \\
Warfarin + Quinidine & Moderate & Increased serum concentration \\
Warfarin + Propranolol & Moderate & Increased serum concentration \\
Amiodarone + Digoxin & Moderate & Increased Digoxin concentration \\
Amiodarone + Quinidine & Moderate & Increased Quinidine concentration \\
Amiodarone + Propranolol & Moderate & Increased therapeutic efficacy \\
Digoxin + Quinidine & Moderate & Increased Digoxin concentration \\
Digoxin + Propranolol & Moderate & Increased arrhythmogenic activity \\
Quinidine + Propranolol & Moderate & Increased Propranolol concentration \\
\bottomrule
\end{tabular}
\end{table}

\subsection{Alternative Therapy Evaluation}

Given Warfarin's high interaction count (4,145 total) and contribution to the only major interaction in this regimen, we evaluated Dabigatran etexilate as a potential therapeutic alternative. Dabigatran etexilate is a direct oral anticoagulant (DOAC) present in the knowledge graph with 2,900 documented interactions.

Examination of Dabigatran etexilate's interactions with the remaining regimen drugs (Amiodarone, Digoxin, Quinidine, Propranolol) revealed four interactions, all classified as Moderate. Specifically, Dabigatran-Amiodarone interaction involves increased Dabigatran serum concentration, representing a severity reduction from the Major Warfarin-Amiodarone interaction.

\subsection{Therapeutic Substitution Outcomes}

Substituting Warfarin with Dabigatran etexilate yielded measurable improvements in the regimen's interaction profile (Table~\ref{tab:improvement}). The number of Major interactions decreased from 1 to 0 (100\% reduction), as the Warfarin-Amiodarone Major interaction was replaced by the Dabigatran-Amiodarone Moderate interaction. The total interaction count remained at 10, with all interactions now classified as Moderate. This represents a meaningful clinical improvement, particularly in reducing bleeding risk associated with the Warfarin-Amiodarone combination.

\begin{table}[H]
\centering
\caption{Improvement Metrics After Warfarin $\rightarrow$ Dabigatran Substitution}
\label{tab:improvement}
\begin{tabular}{lccc}
\toprule
\textbf{Metric} & \textbf{Before} & \textbf{After} & \textbf{Change} \\
\midrule
Major Interactions & 1 & 0 & 100\% reduction \\
Moderate Interactions & 9 & 10 & +1 \\
Total Interactions & 10 & 10 & No change \\
Severity-Weighted Sum & 21 & 20 & 4.8\% reduction \\
\bottomrule
\end{tabular}
\end{table}

The severity-weighted sum decreased from 21 (1 Major $\times$ 3 + 9 Moderate $\times$ 2 = 21) to 20 (10 Moderate $\times$ 2 = 20), representing a 4.8\% improvement in the composite severity metric. While the numerical improvement appears modest, the elimination of the Major bleeding risk interaction represents a clinically meaningful outcome.

% ============================================================================
\section{Figures}
% ============================================================================

\subsection{Polypharmacy Risk Assessment Visualizations}

\begin{figure}[H]
\centering
\includegraphics[width=0.85\textwidth]{figures/polypharmacy_risk_escalation.png}
\caption{Polypharmacy risk escalation showing how overall regimen risk increases with the number of concurrent medications. The non-linear relationship demonstrates that adding each additional drug produces disproportionately larger increases in interaction potential.}
\label{fig:risk_escalation}
\end{figure}

\begin{figure}[H]
\centering
\includegraphics[width=0.85\textwidth]{figures/drug_risk_matrix.png}
\caption{Drug risk matrix displaying interaction severities for common drug combinations. Darker cells indicate more severe interactions, enabling rapid identification of problematic drug pairs.}
\label{fig:risk_matrix}
\end{figure}

\begin{figure}[H]
\centering
\includegraphics[width=0.85\textwidth]{figures/severity_heatmap.png}
\caption{Severity heatmap showing the distribution of interaction severities across drug classes. This visualization reveals systematic patterns in which drug categories most frequently produce severe interactions.}
\label{fig:severity_heatmap}
\end{figure}

\subsection{Alternative Therapy Recommendation Visualizations}

\begin{figure}[H]
\centering
\includegraphics[width=0.9\textwidth]{figures/drug_alternatives_heatmap.png}
\caption{Drug alternatives heatmap showing recommendation scores for various therapeutic substitution options. Higher scores (warmer colors) indicate alternatives that better balance therapeutic equivalence with safety improvement.}
\label{fig:alternatives_heatmap}
\end{figure}

\begin{figure}[H]
\centering
\includegraphics[width=0.9\textwidth]{figures/drug_substitution_network.png}
\caption{Drug substitution network visualization showing recommended alternatives and their therapeutic relationships. Edges connect original drugs to potential substitutes, with edge weights reflecting recommendation strength.}
\label{fig:substitution_network}
\end{figure}

\begin{figure}[H]
\centering
\includegraphics[width=0.9\textwidth]{figures/network_safe_alternatives.png}
\caption{Safe alternatives network highlighting drugs with improved safety profiles compared to high-risk options. This network emphasizes therapeutic clusters where safer substitutions are available.}
\label{fig:safe_alternatives}
\end{figure}

\begin{figure}[H]
\centering
\includegraphics[width=0.9\textwidth]{figures/class_alternatives_summary.png}
\caption{Class-level alternatives summary showing therapeutic substitution patterns across drug categories. This aggregate view supports formulary-level decision making for medication optimization.}
\label{fig:class_alternatives}
\end{figure}

% ============================================================================
\section{Conclusion}
% ============================================================================

The integrated framework combining the Polypharmacy Risk Index with multi-objective drug recommendation provides a systematic, evidence-based approach for managing drug-drug interactions in complex medication regimens. Built upon a comprehensive knowledge graph containing 759,774 drug-drug interactions among 4,313 drugs, the system leverages network-based metrics to quantify risk and identify safer therapeutic alternatives.

The cardiovascular case study demonstrates practical utility of the framework. Substituting Warfarin with Dabigatran etexilate eliminated the Major bleeding risk interaction while maintaining therapeutic coverage for anticoagulation. Although the total number of interactions remained unchanged, the severity profile improved through elimination of the most clinically concerning interaction.

These results suggest that systematic application of the PRI and recommendation system can identify opportunities for medication optimization in polypharmacy scenarios. The framework provides quantified risk metrics and evidence-based alternative suggestions to support clinical decision-making, with all recommendations traceable to documented interactions in the underlying knowledge graph.

% ============================================================================
\section*{References}
% ============================================================================

\bibliographystyle{plain}
\begin{thebibliography}{4}
\bibitem{wishart2018drugbank} Wishart DS, et al. DrugBank 5.0: a major update to the DrugBank database for 2018. \textit{Nucleic Acids Res}. 2018;46(D1):D1074-D1082.
\bibitem{kuhn2016sider} Kuhn M, et al. The SIDER database of drugs and side effects. \textit{Nucleic Acids Res}. 2016;44(D1):D1075-D1079.
\bibitem{davis2021ctd} Davis AP, et al. Comparative Toxicogenomics Database (CTD): update 2021. \textit{Nucleic Acids Res}. 2021;49(D1):D1138-D1143.
\bibitem{beers2019} American Geriatrics Society 2019 Beers Criteria Update Expert Panel. American Geriatrics Society 2019 Updated AGS Beers Criteria. \textit{J Am Geriatr Soc}. 2019;67(4):674-694.
\end{thebibliography}

\end{document}
